% Pregunta 712 (2007-julio-OEP2007, n.º 20): modelo de Barro-Gordon, compromiso frente a discrecionalidad.
% Función de reacción con discrecionalidad: pi = 1,5 + 0,5 pi^e (pendiente b^2/(a+b^2) < 1, ordenada b(k-1)ybar/(a+b^2) > 0).
% Corte con la bisectriz pi = pi^e: D en (3,3). Compromiso: C en (0,0). Tentación con pi^e = 0: T en (0,1,5).
% Autor: Víctor Gutiérrez Marcos
\documentclass[tikz,border=6pt]{standalone}
\input{tikz-tcee}
\begin{document}
\begin{tikzpicture}[tcee,x=1.1cm,y=1.1cm]
  \ejes{$\pi^e$}{$\pi$}{5.2}{4.9}
  % Bisectriz: expectativas racionales (sin sorpresas, y = ybar)
  \draw[curva2] (0,0) -- (4.5,4.5) node[etiqueta,above,text=tceeazul,align=center] {$\pi=\pi^e$\\(sin sorpresa: $y=\bar y$)};
  % Función de reacción del banco central con discrecionalidad
  \draw[curva1] (0,1.5) -- (5.0,4.0) node[etiqueta,right,text=tceemorado,align=left] {reacción del\\banco central};
  % Equilibrio discrecional
  \proyeccion{(3,3)}{$\pi^e_D$}{$\pi_D$}
  \node[punto] at (3,3) {};
  \node[etiqueta,below right] at (3,3) {$D$};
  % Compromiso
  \node[punto] at (0,0) {};
  \node[etiqueta,left] at (-0.08,0.3) {$C$};
  % Tentación: si pi^e = 0, el banco elige pi = 1,5 > pi^e
  \node[circle,draw=tceerojo,fill=white,thick,inner sep=2pt] at (0,1.5) {};
  \node[etiqueta,left,text=tceerojo] at (-0.08,1.55) {$T$};
  \node[etiqueta,right,text=tceerojo,font=\footnotesize,align=left] at (1.5,0.72) {tentación:\\con $\pi^e=0$,\\sorprender\\($y>\bar y$)};
  \draw[flecha,tceerojo] (0.25,0.15) -- (0.25,1.3);
  % Sesgo inflacionista
  \draw[flecha,tceerojo,<->] (-1.0,0) -- node[left,font=\footnotesize,text=tceerojo,align=right] {sesgo\\inflacionista} (-1.0,3);
\end{tikzpicture}
\end{document}
